\documentclass[10pt,aspectratio=43,mathserif]{beamer}		
%设置为 Beamer 文档类型，设置字体为 10pt，长宽比为16:9，数学字体为 serif 风格

%%%%-----导入宏包-----%%%%
\usepackage{ecnu}			%导入 CCNU 模板宏包
\usepackage{ctex}			 %导入 ctex 宏包，添加中文支持
\usepackage{amsmath,amsfonts,amssymb,bm}   %导入数学公式所需宏包
\usepackage{color}			 %字体颜色支持
\usepackage{graphicx,hyperref,url}	
\usepackage{minted,fontspec}
%%%%%%%%%%%%%%%%%%	

\setmonofont{Source Code Pro}
\definecolor{bg}{rgb}{0.95,0.95,0.95}
\newminted{asm}{bgcolor=bg}
\newminted{cpp}{bgcolor=bg}


\beamertemplateballitem		%设置 Beamer 主题

%%%%------------------------%%%%%
\catcode`\。=\active         %或者=13
\newcommand{。}{．}				
%将正文中的“。”号转换为“.”。
%%%%%%%%%%%%%%%%%%%%%

%%%%----首页信息设置----%%%%
\title[ECNUBeamerTemplate]{并行计算大作业}
\subtitle{——并行优化矩阵乘法}			
%%%%----标题设置


\author[x]{
%   ~ \\\medskip
  {\small {朱桐~10175102111}} \\
  {\small {周亦然~10175102207}}
  }
%%%%----个人信息设置

% \institute[ECNU]{
%   East China Normal University \\
%   Department of Computer Science and Technology}
%%%%----机构信息

\date[2019.12.27]{
  2019.12.27}
%%%%----日期信息

\begin{document}

\begin{frame}
\titlepage
\end{frame}				%生成标题页

\section{Outline}
\begin{frame}
\frametitle{Outline}
\tableofcontents
\end{frame}				%生成提纲页

% fastio
\begin{frame}[fragile]
    \frametitle{code: buffered input}
    \begin{cppcode}
    inline char next_char()
    {
        static char buf[100000], *p1 = buf,
            *p2 = buf;
        return p1 == p2 && (p2 = (p1 = buf)
            + fread(buf, 1, 100000, fin), p1 == p2) 
            ? EOF : *p1++;
    }
    \end{cppcode}
\end{frame}


\begin{frame}[fragile]
    \frametitle{code: buffered output}
    \begin{cppcode}
    inline void flush() {
        fwrite(buffer, 1, s-buffer, stdout);
        s = buffer;
        fflush(stdout);
    }
    inline void print(const char ch) {
        // putchar(ch); return;
        if (s-buffer>OutputBufferSize-2) flush();
        *s++ = ch;
    }
    \end{cppcode}
\end{frame}
 
 

\begin{frame}[fragile]
    \frametitle{code: share memory, block partition}
    \begin{cppcode}
    __shared__ ld c_a[max_shared_size];
    int index = blockDim.x * blockIdx.x + threadIdx.x;
    if (index >= an * bm) return;
    int st = min(index, addi) * (workload+1) + 
    max(0, index - addi) * workload, ed = 
    st + workload + (index < addi ? 1 : 0);
    int shareda = min(am, max_shared_size);
    for (int p=st; p<ed; ++p) {
        // ...
    \end{cppcode}
\end{frame} 

\begin{frame}[fragile]
    \frametitle{code: share memory, block partition}
    \begin{cppcode}
    for (int p=st; p<ed; ++p) {
        int i = p / bm, j = p % bm;
        if (p % bm == 0) {
            
            for (int j=0; j<shareda; ++j) {
                c_a[j] = d_a[i * am + j];
            }
            __syncthreads();
        }
        ...
    \end{cppcode}
\end{frame}

\begin{frame}[fragile]
    \frametitle{code: share memory, block partition}
    \begin{cppcode}
    for (int p=st; p<ed; ++p) {
        int i = p / bm, j = p % bm;
        if (p % bm == 0) {
            
            for (int j=0; j<shareda; ++j) {
                c_a[j] = d_a[i * am + j];
            }
            __syncthreads();
        }
        ...
    \end{cppcode}
\end{frame}


\begin{frame}[fragile]
    \frametitle{code: async IO, creating streams}
    \begin{cppcode}
    int st = 0, ed = n * m;
    // printf("st=%d ed=%d, a=%p\n", st, ed, a); 
    cudaStream_t stream[2];
    int mask = 0;
    cudaStreamCreate(&stream[0]);
    cudaStreamCreate(&stream[1]);
    int size;
    \end{cppcode}
\end{frame}

\begin{frame}[fragile]
    \frametitle{code: async IO, creating streams}
    \begin{cppcode}
    cudaStream_t mainstream;
    cudaStreamCreate(&mainstream);
    // ...
    copyMatrixAsync(h_a, d_a, an, am, mainstream);
    // ... read h_b
    copyMatrixAsync(h_b, d_b, bn, bm, mainstream);
    // ... 
    matrixMult<<<grids, block_size, 0, mainstream>>>
    (d_a, d_b, d_c, an, bm, am);
    // ...
    handleCudaError(cudaStreamSynchronize(mainstream));
    // ...
    outputMatrixAsync(h_c, d_c, n, m);
    // pipeline: memcpy from device & output each rows 
    \end{cppcode}
\end{frame}

\begin{frame}[fragile]
    \frametitle{code: async IO, pipelines}
    \begin{cppcode}
    for (; st<ed; st+=size, mask^=1) {
        size = min(chunk_size, ed - st);
        handleCudaError(cudaMemcpyAsync(a + st,
         d_a + st, size * sizeof(ld), 
         cudaMemcpyDeviceToHost, stream[mask]));
        // exit(0);
        if (st - chunk_size >= 0) {
            // printf("%d %d\n",st-chunk_size, st);
            handleCudaError(cudaStreamSynchronize(stream[mask^1]));
            outputinterval(a, st-chunk_size, st);
        }
    }
    \end{cppcode}
\end{frame}



\section{Introduction}
\begin{frame}
  \frametitle{Backgrounds}
  \begin{itemize}
    \item {1111111111111}
    \item {2222222222222}
    \item {3333333333333}
    \item {4444444444444}
  \end{itemize}
\end{frame}

\section{Photo}
\begin{frame}
  \frametitle{My Photo}
    \begin{figure}[htbp]
        \centering
        \includegraphics[width=0.6\textwidth]{pic.png}
        \caption{hahahaha...}\label{fig:digit}
    \end{figure}
\end{frame}

\section{Conclusion}
\begin{frame}
  \frametitle{Emmm...}
  \noindent\textbf{Sequence Tagging Loss\\}
    \[{\mathcal{L}_p} =  - \sum\limits_{i = 1}^S {\sum\limits_{j = 1}^N {{p_{i,j}}\log ({{\hat p}_{i,j}})} } \]
  \noindent\textbf{Language Classifier Loss\\}
   \[{\mathcal{L}_a} =  - \sum\limits_{i = 1}^S {{l_i}\log ({{\hat l}_i})}\]
  \noindent\textbf{Bidirectional Language Model Loss\\}
   \[{\mathcal{L}_l} =  - \sum\limits_{i = 1}^S {\sum\limits_{j = 1}^N {\log (P({w_{j + 1}}|{f_j})) + \log (P({w_{j - 1}}|{b_j}))} }\]
\end{frame}
 
\section{References}
\begin{frame}{References}
\begin{thebibliography}{99}
\bibitem{zhao1} Xuezhe Ma and Eduard Hovy. (2016).\\
{\bf End-to-end Sequence Labeling via Bi-directional LSTM-CNNs-CRF.\\}
In Proceedings of the 54th Annual Meeting of the Association for Computational Linguistics, pages 1064–1074, Berlin, Germany, August 7-12, 2016.
\bibitem{zhao2} Marek Rei. (2017).\\
{\bf Semi-supervised Multitask Learning for Sequence Labeling.\\}
In Proceedings of the 55th Annual Meeting of the Association for Computational Linguistics, pages 2121–2130, Vancouver, Canada, July 30 - August 4, 2017.
\end{thebibliography}
\end{frame}

\end{document}
